% Insertar bajo 2.8 Pruebas; requiere longtable,array,booktabs,graphicx,float,hyperref.
Las pruebas verifican los cuatro requisitos prioritarios: autenticación
(RF-01), pacientes (RF-02), solicitudes radiográficas (RF-03) y acceso
temporal por enlace o QR (RF-04). Se utilizan caminos correctos y errores,
con datos sintéticos, comprobando el estado HTTP y las reglas de acceso.

\subsubsection{Entorno y ejecución}

Las pruebas automatizadas ejecutan el backend real de Express y sus
migraciones SQL sobre PostgreSQL embebido mediante PGlite. Supertest realiza
solicitudes HTTP y Playwright verifica el acceso desde el navegador y genera
evidencias. Este entorno está aislado de los pacientes y archivos reales.
La prueba local de archivos utiliza la ruta multipart y contenido en la base
de prueba; no acredita la carga directa de producción en Supabase Storage.

La ejecución del 11/10/2026 a las 22:20:56 (hora de Bolivia)
registró 31 pruebas, 31 aprobadas y 0 fallidas, sin omitidas.
El escenario adicional de navegador se aprobó y verificó 15
respuestas HTTP. El archivo de seguridad específico registra 22
comprobaciones HTTP; estos conteos describen distintos niveles y no deben sumarse
como si fueran casos independientes.

Los comandos reproducibles son \texttt{npm run test:report} y
\texttt{npm run test:e2e -w backend -- security.spec.js}.
El reporte completo se conserva en
\path{evidencia/reporte-pruebas-e3.txt}; las respuestas resumidas se encuentran
en \path{evidencia/seguridad-http.json} y
\path{evidencia/evidencias-navegador.json}. Ninguno incluye contraseñas ni JWT.

\subsubsection{Casos ejecutados}

\begingroup
\footnotesize
\setlength{\tabcolsep}{3pt}
\renewcommand{\arraystretch}{1.2}
\begin{longtable}{@{}>{\raggedright\arraybackslash}p{.09\linewidth}>{\raggedright\arraybackslash}p{.24\linewidth}>{\raggedright\arraybackslash}p{.21\linewidth}>{\raggedright\arraybackslash}p{.27\linewidth}>{\raggedright\arraybackslash}p{.11\linewidth}@{}}
\caption{Casos de prueba ejecutados y resultados observados}\label{tab:pruebas-reales}\\
\toprule
\textbf{ID}&\textbf{Escenario}&\textbf{Esperado}&\textbf{Obtenido y evidencia}&\textbf{Estado}\\\midrule
\endfirsthead
\toprule
\textbf{ID}&\textbf{Escenario}&\textbf{Esperado}&\textbf{Obtenido y evidencia}&\textbf{Estado}\\\midrule
\endhead
\bottomrule\endlastfoot
CP-01 & RF-01: credenciales válidas & HTTP 200 y JWT & HTTP 200; EV-01; Figura~\ref{fig:prueba-200} & Aprobado \\
CP-02 & RF-01: contraseña incorrecta & HTTP 401, mensaje genérico & HTTP 401; EV-02 & Aprobado \\
CP-03 & Sin JWT en pacientes & HTTP 401 & HTTP 401; EV-03; Figura~\ref{fig:prueba-401} & Aprobado \\
CP-04 & Odontólogo consulta cuentas & HTTP 403 & HTTP 403; EV-08; Figura~\ref{fig:prueba-403} & Aprobado \\
CP-05 & RF-02: paciente sin campos obligatorios & HTTP 422 & HTTP 422; EV-04; Figura~\ref{fig:prueba-422} & Aprobado \\
CP-06 & RF-02: crear, consultar, actualizar y eliminar paciente & 201, 200, 200, 200; posterior 404 & Códigos coincidentes; SEG-04 a 06 y SEG-21 a 22 & Aprobado \\
CP-07 & Paciente ajeno al odontólogo & HTTP 404 & HTTP 404; SEG-11 & Aprobado \\
CP-08 & RF-03: crear solicitud de paciente asignado & HTTP 201 y asociación persistida & HTTP 201; EV-11; verificación SQL en radiography.test.js & Aprobado \\
CP-09 & RF-03: solicitar estudio sin asignación & HTTP 404 & HTTP 404; EV-09 & Aprobado \\
CP-10 & RF-04: generar enlace y QR & HTTP 201, URL y PNG del QR & HTTP 201; EV-12 y EV-13 & Aprobado \\
CP-11 & RF-04: validar acceso vigente & HTTP 200, datos mínimos & HTTP 200; EV-15 & Aprobado \\
CP-12 & RF-04: reutilizar enlace revocado & HTTP 410 & HTTP 410; EV-14 & Aprobado \\
CP-13 & Enlace vencido y carga repetida & HTTP 410 en ambos casos & HTTP 410; radiography.test.js & Aprobado \\
CP-14 & Carga multipart local y consulta protegida & HTTP 201 y 200; solicitud completada & Códigos coincidentes; radiography.test.js & Aprobado \\
CP-15 & Archivo con contenido o extensión inválidos; exceso de tamaño & 422, 422 y 413 & Códigos coincidentes; radiography.test.js & Aprobado \\
CP-16 & Paciente intenta ingresar con contraseña válida o JWT anterior & HTTP 401; sin token nuevo & HTTP 401; SEG-16 y SEG-17 & Aprobado \\
CP-17 & JWT alterado, vencido o reutilizado después de logout & HTTP 401 & HTTP 401; SEG-13, SEG-18 y SEG-20 & Aprobado \\
CP-18 & Odontólogo usa otra clínica & HTTP 403 & HTTP 403; SEG-14 & Aprobado \\
CP-19 & Permiso directo de cuentas concedido a odontólogo & HTTP 403 también en ruta anterior & HTTP 403; SEG-15 & Aprobado \\
CP-20 & Recepcionista intenta acceder a radiografías y datos clínicos & HTTP 403 y omisión de campos sensibles & Comprobado; receptionist.test.js & Aprobado \\
\end{longtable}
\noindent\textit{Fuente: elaboración propia a partir de la ejecución automatizada.}
\endgroup

Los identificadores EV corresponden al informe de navegador y los SEG al
informe HTTP de seguridad. Los archivos de pruebas citados pertenecen a
\path{backend/test/}. Las filas agrupan operaciones relacionadas; el número
de filas no equivale al contador del ejecutor de Node.js.

\subsubsection{Evidencias}

Las siguientes capturas muestran respuestas reales del servidor de prueba
registradas por Playwright. Son capturas del informe de ejecución, con las
credenciales y los tokens omitidos, y no capturas del despliegue de Vercel.

\begin{figure}[H]
\centering
\includegraphics[width=\textwidth]{../../evidencia/http-200.png}
\caption{Inicio de sesión autorizado, CP-01. Fuente: ejecución propia en entorno local.}
\label{fig:prueba-200}
\end{figure}

\begin{figure}[H]
\centering
\includegraphics[width=\textwidth]{../../evidencia/http-401.png}
\caption{Acceso a pacientes sin JWT, CP-03. Fuente: ejecución propia en entorno local.}
\label{fig:prueba-401}
\end{figure}

\begin{figure}[H]
\centering
\includegraphics[width=\textwidth]{../../evidencia/http-403.png}
\caption{Consulta de cuentas por un odontólogo, CP-04. Fuente: ejecución propia en entorno local.}
\label{fig:prueba-403}
\end{figure}

\begin{figure}[H]
\centering
\includegraphics[width=\textwidth]{../../evidencia/http-422.png}
\caption{Registro de paciente sin campos obligatorios, CP-05. Fuente: ejecución propia en entorno local.}
\label{fig:prueba-422}
\end{figure}

El informe ampliado \path{evidencia/evidencias-seguridad.pdf} contiene las
15 respuestas del recorrido de navegador, incluidas la creación de la
solicitud, la generación del enlace y el rechazo del enlace revocado.

\subsubsection{Comprobaciones en producción y rendimiento}

En el despliegue de Vercel se realizaron comprobaciones de lectura:
la salud de la API respondió 200, el acceso a pacientes sin autenticación
respondió 401, el inicio de sesión respondió 200 y la consulta de solicitudes
radiográficas respondió 200. El detalle se conserva en
\path{evidencia/verificacion-despliegue.json}.

Se midieron diez consultas consecutivas autenticadas al listado de pacientes,
con página de quince registros y sin transferencias de archivos. La medición
incluye el tiempo observado por el cliente HTTP desde este equipo, la red y
la lectura de la respuesta JSON. No es una prueba de concurrencia ni de carga.

\begin{table}[H]\centering
\caption{Tiempos observados en consultas de pacientes en Vercel}
\begin{tabular}{rrr}\toprule
\textbf{Consulta}&\textbf{Estado HTTP}&\textbf{Tiempo (ms)}\\\midrule
1&200&349\\
2&200&214\\
3&200&209\\
4&200&251\\
5&200&219\\
6&200&222\\
7&200&208\\
8&200&320\\
9&200&213\\
10&200&217\\
\bottomrule\end{tabular}\end{table}
En esta muestra, 10 de 10 consultas respondieron en menos de tres
segundos, con un rango de 208 a 349 ms.
Se satisface el umbral de RNF-01 para esta operación y estas condiciones;
el resultado no se generaliza a todas las operaciones ni a usuarios concurrentes.

Las modificaciones de autorización y las nuevas rutas v1 documentadas en esta
revisión se verificaron localmente. El recorrido completo de creación,
enlace y entrega directa a Supabase Storage no se ejecutó en producción
durante esta revisión; las comprobaciones de lectura anteriores no lo
sustituyen. Antes de presentar la versión desplegada debe verificarse ese
recorrido sobre la misma revisión del código.
